\documentclass[11pt,a4paper]{article}
\usepackage[T1]{fontenc}
\usepackage[margin=25mm,headheight=15pt]{geometry}
\usepackage{fancyhdr}
\usepackage{multicol}
\usepackage{pdflscape}
\usepackage{booktabs}
\pagestyle{fancy}
\fancyhf{}
\lhead{Technical Writing}
\rhead{Page Layout}
\cfoot{\thepage}

\title{Page Layout in LaTeX}
\author{Your Name}
\date{\today}

\begin{document}
\maketitle
\section{Paper size and margins}
This document uses A4 paper and 25 mm margins. Change the options in the geometry package to experiment with other page sizes and margins.

\section{Headers and footers}
The running header displays the course and topic. The footer displays the page number.

\section{Multiple columns}
\begin{multicols}{2}
Technical documents sometimes use columns for compact reference material. A two-column layout should remain readable and should not force the reader to jump between unrelated ideas.

This second paragraph continues the example in the same column layout. Try changing the number of columns and observe how the text is balanced.
\end{multicols}

\begin{landscape}
\section{Landscape page}
This page is in landscape orientation. It can be used for content that is wider than a portrait page.
\begin{center}
\begin{tabular}{lll}
\toprule
Setting & Purpose & Example\\
\midrule
Paper size & Page dimensions & A4\\
Margins & Writing area & 25 mm\\
Orientation & Page direction & Landscape\\
Columns & Reading layout & Two columns\\
\bottomrule
\end{tabular}
\end{center}
\end{landscape}
\end{document}
